\subsection{The category of sheaves}

As discussed, presheaves form a category for which sheaves form a full subcategory.
We will denote the category of $\C$-valued presheaves over a space $X$ by $\C_X^\pre$ (instead of $\Psh_\C(X)$), and we will denote the full
subcategory of sheaves by $\C_X\subseteq\C_X^\pre$.

Note that if $\F,G\in\C_X$ are sheaves (or presheaves), then the hom-set of sheaf morphisms $\F\to\G$ is indeed a set, since $X$ as a category
is small (functor categories under small categories are locally small).
Furthermore, if $\C$ is enriched over some monoidal category (in our case, usually $\Ab$) so is $\C_X$.
This is because it is true in general for functor categories over enriched categories.

\bdefn

    Let $\F,\G$ be $\C$-valued sheaves defined on $X$.
    Define their {\emphcolor sheaf-hom} to be the sheaf whose sections are sheaf morphisms,
    $$ \ihom(\F,\G)(U) = \hom_{\C_X}(\F|_U,\G|_U) . $$
    As discussed, if $\C$ is enriched over a monoidal category $\V$ then so is $\C_X$, and so $\ihom(\F,\G)$ takes values in $\V$.
    In particular, if $\F,\G$ are sheaves of Abelian groups, then so is their sheaf-hom.

\edefn

\bprop

    The sheaf-hom is indeed a sheaf.

\eprop

\bproof

    We first show it is a sheaf: if $U\subseteq V$ then a sheaf morphism $\F|_V\to\G|_V$ restricts to a sheaf morphism
    $\F|_U\to\G|_U$ in the obvious way: its components are the same but only taken on subsets of $U$.
    This is clearly a presheaf.

    To show that it is a sheaf, let $\set{U_i}$ be an open cover for $U\subseteq X$, and $\alpha_i\colon\F|_{U_i}\to\G|_{U_i}$ be sheaf morphisms
    such that $\alpha_i|_{U_i\cap U_j}=\alpha_j|_{U_i\cap U_j}$ for all $i,j$.
    In order to glue them together for a sheaf morphism $\alpha\colon\F_U\to\G_U$ we need to define $\alpha_W$ for all $W\subseteq U$.
    For $s\in\F(W)$ consider $s_i=s|_{U_i}\in U_i\cap W$, and its image under $\alpha_i$ to $\G(U_i\cap W)$, call it $t_i$.
    Then since $U_i\cap W$ cover $W$ and since $\alpha$ is a morphism as well as $\G$ being a sheaf, we can uniquely glue $t_i$ together to
    get a $t\in\G(W)$; this will be the image of $s$.

    It is not hard to see that this is the unique gluing of $\alpha_i$, and is itself a sheaf morphism.
    \qed

\eproof

\bnote

    In the above proof we needed only for $\G$ to be a sheaf, $\F$ can be a presheaf.

    Furthermore $\ihom$ forms a bifunctor
    $$ \ihom\colon\C_X^\op\times\C_X \to \Set_X , $$
    or $V_X$ instead of $\Set_X$ when $\C$ is enriched over $V$.

\enote

Note that if $\O_X$ is a ringed space, it is also a module over itself.
Thus if $\F$ is an $\O_X$ module, we can form its {\emphcolor dual}:
$$ \F^\vee = \ihom_{\Mod_{\O_X}}(\F,\O_X) . $$

We now focus on sheaves of Abelian groups.
It is a classical result that functor categories over Abelian categories are Abelian, thus $\Ab_X^\pre$ is an Abelian category.
That is to say, the hom-sets (not hom-sheaves) have an Abelian group structure: for $\phi,\psi\colon\F\to\G$ we can define $\phi+\psi$
and $-\phi$ by $(\phi+\psi)(U)=\phi(U)+\psi(U)$ and $(-\phi)(U)=-\phi(U)$).
And we can form kernels and cokernels so that all monomorphisms are kernels and all epimorphisms are cokernels.

As with all limits in functor categories, these are done coordinate-wise.
That is to say, if we denote by $\ker_\pre$ and $\coker_\pre$ the kernel and cokernel of Abelian presheaf morphisms, we have
$$ (\ker_\pre\phi)(U) = \ker(\phi(U)),\qquad (\coker_\pre\phi)(U) = \coker(\phi(U)) . $$

\bremark[name={monepiccomp}]

    Let $I$ be an indexing category and $\C$ be a category with arbitrary pullbacks.
    It is easy to see that a $\C$-morphism $f\colon X\to Y$ is monic iff the following square is a pullback:

    \commsq{X&X\cr X&Y}
        {\id}{f}{\id}{f}

    Because the functor category $[I,\C]$ also has arbitrary pullbacks, a natural transformation $\alpha\colon\F\To\G$ is monic iff
    a similar diagram is a pullback.
    Because limits in a functor category are taken coordinate-wise, this is iff $\alpha(i)$ is monic for all $i\in I$.

    Thus a natural transformation is monic iff its coordinates are.

    Similarly, if $\C$ has arbitrary pushouts, natural transformations in $[I,\C]$ are epic iff their coordinates are.

\eremark

In lieu of this note, we see that $\Ab^\pre_X$ forms an Abelian category.
Indeed, if $\phi\colon\F\To\G$ is a monic morphism of presheaves, then $\phi(U)\colon\F(U)\To\G(U)$ is a monic morphism of Abelian groups
for all $U\subseteq X$.
Thus $\phi(U)=\ker\coker\phi(U)$ for all $U\subseteq X$ and thus $\phi=\ker\coker\phi$ as desired.
Similar for epic morphisms.

Many properties of $\Ab^\pre_X$ are seen to be pointwise.
For one, if we set an open set $U\subseteq X$, the functor $\Ab^\pre_X\to\Ab$ mapping $\F\mapsto\F(U)$ is an exact functor.
Indeed, it is additive and preserves kernels and cokernels because these are taken pointwise.

Furthermore, a sequence of presheaves
$$ \F_1\to\F_2\to\cdots\to\F_n $$
is exact iff $\F_1(U)\to\cdots\to\F_n(U)$ is for all $U\subseteq X$.
Again, this is an immediate consequence of kernels and cokernels (and thus images) being taken coordinate-wise.

Everything we have said holds for presheaves valued in any Abelian category $\A$ (recall that Abelian categories have all finite limits and colimits).
But we are more interested in the category of sheaves, in which case this is less obvious.

Note that the constant presheaf $\ul0_\pre$ is a sheaf, so we can denote it by $\ul0$ without confusion.
In order to show that $\A_X$ is an Abelian category, we must first show that kernels and cokernels exist in it.
The following proposition shows that the kernel of a sheaf morphism is the same as the kernel of a presheaf morphism, and thus we get the existence
of kernels for free.

\bprop

    If $\phi\colon\F\to\G$ is a morphism of sheaves, then $\ker_\pre\phi$ is a sheaf.
    Thus $\ker\phi=\ker_\pre\phi$ forms a kernel in the category of sheaves.

\eprop

\bproof

    We show this for $\A=\Ab$: if $\set{U_i}_i$ is an open cover of $U\subseteq X$ and $s_i\in\ker\phi(U_i)$ are such that
    $s_i|_{U_i\cap U_j}=s_j|_{U_i\cap U_j}$ for all $i,j$, then by gluing they form a section $s\in\F(U)$.
    And by naturality
    $$ \phi(U)(s)|_{U_i}=\phi(U_i)(s|_{U_i})=\phi(U_i)(s_i) = 0 . $$
    The unique gluing of $0\in\G(U_i)$ is $0\in\G(U)$, so that $\phi(U)(s)=0$, meaning $s\in\ker\phi(U)$ as desired.
    \qed

\eproof

\bremark[name={sheafinclim}]

    A generalization of the above proposition shows that limits taken in the category of presheaves of sheaves forms a sheaf.
    That is, if $\Phi\colon I\to\C_X$ is a functor, then consider the inclusion $\iota\colon\C_X\inj\C_X^\pre$.
    Then $\lim\iota\circ\Phi$ is a sheaf, so $\lim\Phi$ exists in $\C_X$.

    Thus if $\C$ has arbitrary pullbacks, so too does $\C_X^\pre$ and thus so too does $\C_X$.
    It then follows that a morphism of sheaves is monic iff its monic as a morphism of presheaves (considering the pullback condition
    \refmath[remark]{monepiccomp}) iff each component is monic (again by \refmath[remark]{monepiccomp}).

\eremark

Unfortunately we now run in to our first hurdle: presheaf cokernels are not sheaves.

\bexam

    Consider $\bC$ with its usual topology, and let $\O_\bC$ be the sheaf of holomorphic functions.
    Let $\F$ be the presheaf of functions admitting a holomorphic logarithm.
    Consider the following exact sequence of presheaves
    $$ 0 \to \ul\bZ \inj \O_\bC \xtos{\exp(2\pi i)} \F \to 0 , $$
    where the second map is $f\mapsto\exp(2\pi if)$.
    This is indeed an exact sequence of presheaves: the first morphism is an embedding, and the second is epic (if $f=\exp(g)$ then
    $g/2\pi i\in\O_\bC$ is an origin for $f$).
    The composition is zero, as $\exp(2\pi in)=0$, and if $\exp(2\pi if)=0$ then $f(x)\in\bZ$ for all $x\in\bC$.
    Because $\bZ$ is discrete, this means $f\in\ul\bZ$ as desired.

    This shows that $\F$ is the cokernel of $\ul\bZ\inj\O_\bC$: but it is not a sheaf!
    Indeed $\id\colon\bC-0\to\bC$ has a local logarithm at every point in its domain, but no global logarithm.
    Thus $\F$ fails the gluing axiom.

\eexam

Thus for now we must wait to demonstrate that $\A_X$ is an Abelian category.

